% Moderne Seitenrand-Konfiguration für LuaLaTeX
% Optimiert für sowohl Druck als auch digitale Ansicht

\usepackage{geometry}

% --- Automatische Erkennung: Druck vs. Digital ---
% Option 1: Manuelle Wahl (kommentieren Sie eine Zeile aus)
\newif\ifdigitalversion
\digitalversiontrue    % Für PDF/Digital
% \digitalversionfalse  % Für Druck

% --- Seitenränder basierend auf Version ---
\ifdigitalversion
  % === DIGITALE VERSION (PDF) ===
  \geometry{
    a4paper,
    twoside,                    % Zweiseitig auch für PDF (bessere Navigation)
    left=30mm,                  % Symmetrische Ränder
    right=30mm,
    top=25mm,
    bottom=30mm,
    % Moderne Ergänzungen
    headheight=15pt,            % Genug Platz für Kopfzeile
    headsep=10mm,               % Abstand Kopfzeile-Text
    footskip=12mm,              % Abstand Text-Fußzeile
    marginparwidth=20mm,        % Randnotizen (falls verwendet)
    marginparsep=5mm,           % Abstand zu Randnotizen
    % PDF-optimiert
    includeheadfoot=false,      % Kopf/Fuß außerhalb Satzspiegel
    bindingoffset=0mm           % Keine Bindekorrektur für PDF
  }
\else
  % === DRUCK-VERSION ===
  \geometry{
    a4paper,
    twoside,
    left=35mm,                  % Mehr Platz für Bindung
    right=25mm,                 % Weniger rechts (Daumengriff)
    top=25mm,
    bottom=30mm,
    % Druck-spezifisch
    headheight=15pt,
    headsep=10mm,
    footskip=12mm,
    marginparwidth=15mm,        % Weniger Randnotizen-Platz
    marginparsep=3mm,
    includeheadfoot=false,
    bindingoffset=5mm           % Zusätzlich für Bindung
  }
\fi

% --- Alternative: Umgebungsbasierte Erkennung ---
% Automatische Erkennung basierend auf Compile-Flags
% Verwenden Sie: lualatex -jobname=meinedatei-digital main.tex
% \ifdefstring{\jobname}{.*-digital.*}{
%   \digitalversiontrue
% }{
%   \digitalversionfalse
% }

% --- Debug-Ausgabe ---
\makeatletter
\AtBeginDocument{%
  \ifdigitalversion
    \typeout{=== DIGITALE VERSION ===}
    \typeout{Ränder: links/rechts symmetrisch (30mm)}
  \else
    \typeout{=== DRUCK-VERSION ===}
    \typeout{Ränder: links 35mm, rechts 25mm + Bindekorrektur}
  \fi
  \typeout{Textbreite: \the\textwidth}
  \typeout{Texthöhe: \the\textheight}
  \typeout{=======================}
}
\makeatother

% --- Zusätzliche Layout-Befehle ---
% Umschalten während des Dokuments (falls nötig)
\newcommand{\setdigitalmode}{%
  \digitalversiontrue
  \newgeometry{
    left=30mm, right=30mm, top=25mm, bottom=30mm,
    headheight=15pt, headsep=10mm, footskip=12mm
  }
}

\newcommand{\setprintmode}{%
  \digitalversionfalse
  \newgeometry{
    left=35mm, right=25mm, top=25mm, bottom=30mm,
    headheight=15pt, headsep=10mm, footskip=12mm,
    bindingoffset=5mm
  }
}

% --- Moderne Ergänzungen ---
% Verhindert zu weite Ränder bei langen Wörtern
\raggedbottom                   % Flexible Seitenhöhe

% Optional: Satzspiegel anzeigen (für Layout-Tests)
% \usepackage{showframe}
% \usepackage{layout}           % \layout Befehl für Seitenanalyse

% --- KOMA-kompatible Alternative (falls bevorzugt) ---
% Statt geometry könnte man auch KOMA-Script's typearea verwenden:
% \KOMAoptions{DIV=12, BCOR=5mm}  % Automatische Berechnung